\section{Introduction}
\label{sec:introduction}

% WRITING PROMPTS (comments only; they do not print). About 1 to 1.5 pages,
% five paragraphs. End with a list of contributions.
%
% 1. Context. Why robotic lawn mowing: labour, regular cutting, campus and
%    park lawns. Most commercial mowers still bounce randomly inside a
%    boundary wire. What does random mowing cost in time, overlap and energy?
% 2. Coverage path planning. Define complete coverage path planning (CCPP) in
%    one sentence. Say how the field splits: cell decomposition (boustrophedon),
%    grid methods, and soft computing (GA, PSO, ACO) to order sub-regions.
% 3. The gap. Ask yourself: what do these methods minimise? Mostly length or
%    the number of turns, measured after planning. For a skid-steer robot a
%    point turn is expensive: with the placeholder robot one 90 degree turn
%    costs as much as 3.3 m of straight driving (Fig. energy-model), and turns
%    are 26 to 31 % of the plan energy. So length is a poor proxy for battery use.
% 4. This paper. The faculty robot: a radio-controlled 4WD skid-steer mower.
%    You make it autonomous (Fig. robot-architecture) and plan its coverage by
%    energy. Name the two levers of the method: the decomposition direction
%    and the cell order.
% 5. Contributions, as a list (keep each to one or two lines):
%    - an energy-aware CCPP that chooses the decomposition direction and the
%      cell order by a skid-steer energy model (GA, with a PSO variant);
%    - a fair comparison with a re-implementation of Xie et al. (2024) on
%      shared cells, lanes and transitions, ten lawns, equal budgets, and
%      paired statistical tests, with the gain split into direction and order;
%    - the retrofit that makes a radio-controlled mower autonomous without
%      losing manual control (controller, PWM relay for the blade, sensors);
%    - open code and data (repository and Zenodo DOI).
% Close with one sentence on the structure of the paper (Sections 2 to 9).
